\section*{Capítulo \ref{ch:b2:meiosis-heredity} --- Meiose, mistura genética e hereditariedade}

\begin{solution}{exo:b2:meiosis-heredity:1}
Mitose: uma divisão, sem pareamento de homólogos, duas células, cada
uma geneticamente idêntica à célula-mãe ($2n \to 2n$). \omterm{def:b2:meiosis-heredity:meiosis}{Meiose}: duas
divisões após uma única replicação, os homólogos pareiam e fazem
crossing-over, quatro células, \omterm{def:b2:meiosis-heredity:meiosis}{haploides} e todas geneticamente
diferentes ($2n \to n$).
\end{solution}

\begin{solution}{exo:b2:meiosis-heredity:2}
$2^{23} = 8.4\times 10^{6}$; $2^{4} = 16$; $2^{7} = 128$ --- antes que
o crossing-over multiplique cada um desses números sem limite.
\end{solution}

\begin{solution}{exo:b2:meiosis-heredity:3}
$Tt\times Tt$: \omterm{def:b2:meiosis-heredity:vocabulary}{genótipos} $1\,TT : 2\,Tt : 1\,tt$, \omterm{def:b2:meiosis-heredity:vocabulary}{fenótipos} $3$
altas : $1$ anã. $Tt\times tt$: $1\,Tt : 1\,tt$, $1 : 1$. Faça o
cruzamento-teste da planta alta com uma anã: todas altas indica $TT$,
metade anãs indica $Tt$ (ou autofecunde-a: $TT$ se mantém pura, $Tt$
segrega $3 : 1$).
\end{solution}

\begin{solution}{exo:b2:meiosis-heredity:4}
Genes ligados estão no mesmo cromossomo, e suas combinações parentais
de \omterm{def:b2:meiosis-heredity:vocabulary}{alelos} são super-representadas entre os gametas; a \omterm{def:b2:meiosis-heredity:linkage}{frequência de
recombinação} é a fração de gametas recombinantes (descendentes
recombinantes de um cruzamento-teste); um \omterm{def:b2:meiosis-heredity:linkage}{centimorgan} corresponde a um
por cento de recombinação. Genes muito afastados num cromossomo longo
têm pelo menos um crossing-over entre eles em quase toda \omterm{def:b2:meiosis-heredity:meiosis}{meiose}, e os
crossing-overs múltiplos tornam as combinações aleatórias, de modo que
$r$ chega a $\tfrac12$ --- indistinguível da \omterm{thm:b2:meiosis-heredity:mixing}{segregação independente}.
\end{solution}

\begin{solution}{exo:b2:meiosis-heredity:5}
Total 556; esperados $312.75 : 104.25 : 104.25 : 34.75$; $\chi^2 =
2.25^2/312.75 + 3.75^2/104.25 + 3.25^2/104.25 + 2.75^2/34.75 = 0.016
+ 0.135 + 0.101 + 0.218 = 0.47 < 7.81$: compatível com $9 : 3 : 3 :
1$.
\end{solution}

\begin{solution}{exo:b2:meiosis-heredity:6}
Contra $1 : 1 : 1 : 1$ (250 cada): $\chi^2 = (162^2 + 138^2 + 147^2
+ 153^2)/250 = 361 \gg 7.81$: ligados. Recombinantes $103 + 97 = 200$
em 1000: $r = 0.20$, \qty{20}{cM}. Em repulsão, $Ab/aB$ dá 400
$Ab$, 400 $aB$, 100 $AB$, 100 $ab$: a mesma distância, com as classes
parentais e recombinantes trocadas.
\end{solution}

\begin{solution}{exo:b2:meiosis-heredity:7}
$d = -\tfrac12\ln(1 - 2r)$: $r = 0.20$ dá $-\tfrac12\ln 0.6 =
0.255$, \qty{25.5}{cM}; $r = 0.45$ dá $-\tfrac12\ln 0.1 = 1.15$,
\qty{115}{cM}. $r = \tfrac12(1 - e^{-2d})$: \qty{30}{cM} dão
$\tfrac12(1 - e^{-0.6}) = 0.226$; \qty{100}{cM} dão $\tfrac12(1 -
e^{-2}) = 0.43$. Abaixo de \qty{10}{cM}, a correção é desprezível;
acima de \qty{30}{cM}, é essencial, e acima de \qty{50}{cM} $r$ satura,
de modo que um único cruzamento de dois pontos já não consegue medir a
distância --- mapeiam-se genes distantes encadeando genes próximos.
\end{solution}

\begin{solution}{exo:b2:meiosis-heredity:8}
(a) $X^wX^w \times X^+Y$: todos os filhos brancos ($X^wY$), todas as
filhas vermelhas ($X^+X^w$). (b) $X^+X^w \times X^wY$: filhos metade
vermelhos, metade brancos; filhas metade vermelhas ($X^+X^w$), metade
brancas ($X^wX^w$). (c) filhas
$X^+X^w$ $\times$ irmãos $X^wY$: como em (b) --- metade de cada sexo
branca, metade vermelha.
\end{solution}

\begin{solution}{exo:b2:meiosis-heredity:9}
Duas enzimas em série produzem o pigmento; a cor púrpura exige um
\omterm{def:b2:meiosis-heredity:vocabulary}{alelo} \omterm{def:b2:meiosis-heredity:vocabulary}{dominante} nos dois locos ($C\_P\_$). As linhagens são $CCpp$ e
$ccPP$ (cada uma branca por uma razão diferente); o híbrido $CcPp$ é
púrpura; autofecundado, dá $9\,C\_P\_$ púrpuras para $7$ brancas
($3\,C\_pp + 3\,ccP\_ +
1\,ccpp$). Cruzamento-teste $CcPp\times ccpp$: $1$ púrpura ($CcPp$) :
$3$ brancas.
\end{solution}

\begin{solution}{exo:b2:meiosis-heredity:10}
Parentais $+++$, $abc$; crossing-overs duplos $+b+$, $a+c$; $b$ é o
gene que mudou: ordem $a$--$b$--$c$. $r_{ab} = (11 + 9 + 2)/1202
= 0.018$ (\qty{1.8}{cM}); $r_{bc} = (118 + 112 + 2)/1202 = 0.193$
(\qty{19.3}{cM}). Duplos esperados $0.018\times 0.193\times 1202 =
4.2$; observados 2: coincidência $0.47$, interferência $0.53$.
\end{solution}

\begin{solution}{exo:b2:meiosis-heredity:11}
Ascos de segunda divisão $300/1000 = \qty{30}{\%}$; distância $=
\tfrac12\times 30 = \qty{15}{cM}$. Um crossing-over entre o gene e o
centrômero envolve apenas duas das quatro cromátides, de modo que um
asco com segregação na segunda divisão contém duas cromátides
recombinantes e duas parentais: a \omterm{def:b2:meiosis-heredity:linkage}{frequência de recombinação} é metade
da frequência desses ascos. Para $2 : 4 : 2$: um crossing-over entre o
loco e o centrômero troca as duas cromátides internas, de modo que
cada metade do fuso na \omterm{def:b2:meiosis-heredity:meiosis}{meiose} I leva uma cromátide $A$ e uma $a$; a
\omterm{def:b2:meiosis-heredity:meiosis}{meiose} II as separa então na ordem $A\,a \mid
a\,A$ (ou a inversa), e a mitose final duplica cada esporo:
$2 : 4 : 2$.
\end{solution}

\begin{solution}{exo:b2:meiosis-heredity:12}
A mãe dela é portadora obrigatória (recebeu o único X do pai), de modo
que a mulher recebeu o \omterm{def:b2:meiosis-heredity:vocabulary}{alelo} com probabilidade $\tfrac12$. Dois filhos
saudáveis: uma portadora tem probabilidade $(\tfrac12)^2 = \tfrac14$
de ter dois filhos saudáveis; uma não portadora, probabilidade 1.
Bayes:
$P(\text{portadora}\mid\text{dados}) = (\tfrac12\times\tfrac14)/(\tfrac12
\times\tfrac14 + \tfrac12\times 1) = \tfrac{1/8}{5/8} = \tfrac15$.
\end{solution}

\begin{solution}{pb:b2:meiosis-heredity:1}
\textbf{1.} Genitores $vg/vg\;+/+$ e $+/+\;e/e$; $F_1$ $+/vg\;+/e$;
gametas $+\,+$, $+\,e$, $vg\,+$, $vg\,e$, cada um $\tfrac14$.
\textbf{2.} $900 : 300 : 300 : 100$.
\textbf{3.} $\chi^2 = 8^2/900 + 9^2/300 + 4^2/300 + 3^2/100 = 0.07 +
0.27 + 0.05 + 0.09 = 0.48 < 7.81$: a \omterm{thm:b2:meiosis-heredity:mixing}{segregação independente} se
confirma.
\textbf{4.} Os \omterm{def:b2:meiosis-heredity:vocabulary}{homozigotos} selvagens nos dois locos são $\tfrac{1}{16}$
do total, e $\tfrac{9}{16}$ são selvagens: $\tfrac19$.
\textbf{5.} Selvagem, vestigial, ébano, vestigial ébano, $\tfrac14$
cada.
\textbf{6.} O \omterm{def:b2:meiosis-heredity:vocabulary}{fenótipo} é produto do \omterm{def:b2:meiosis-heredity:vocabulary}{genótipo} e do ambiente (a
temperatura age sobre as enzimas do pigmento): um \omterm{def:b2:meiosis-heredity:vocabulary}{genótipo} especifica
uma norma de reação, e não um caráter fixo.
\textbf{7.} $b\,+/+\,vg$ (repulsão): $b$ veio com $+_{vg}$ de um
genitor, $vg$ com $+_b$ do outro.
\textbf{8.} Parentais: pretas ($b\,+$) 965 e vestigiais ($+\,vg$) 944;
recombinantes: pretas vestigiais ($b\,vg$) 206 e selvagens ($+\,+$)
185. O cromossomo selvagem $+\,+$ não existia na $F_1$; ele só pode ser
formado por um crossing-over.
\textbf{9.} $r = 391/2300 = 0.17$: \qty{17}{cM}.
\textbf{10.} $d = -\tfrac12\ln(1 - 0.34) = -\tfrac12\ln 0.66 = 0.21$:
\qty{21}{cM}.
\textbf{11.} Não há crossing-over nos machos de \emph{Drosophila}: os
gametas do macho levam apenas os dois cromossomos parentais.
\textbf{12.} $F_1$ $b\,vg/+\,+$ (acoplamento): classes parentais
selvagem e preta vestigial, \qty{41.5}{\%} cada; recombinantes preta
e vestigial, \qty{8.5}{\%} cada.
\textbf{13.} $X^hX^+$: uma filha recebe o único X do pai, que leva $h$,
e um X normal da mãe não afetada.
\textbf{14.} $\tfrac12$ (menino) $\times$ $\tfrac12$ (recebe $X^h$) $=
\tfrac14$.
\textbf{15.} $\tfrac12$: ela recebe um ou outro dos cromossomos X da
mãe; o pai lhe dá um X normal.
\textbf{16.} Cada filho é não afetado com probabilidade $\tfrac34$:
$(\tfrac34)^3 = \tfrac{27}{64} = 0.42$.
\textbf{17.} Um filho recebe o Y do pai, nunca o X. Uma mulher só é
afetada se for $X^hX^h$: pai afetado e mãe portadora (ou afetada), o
que é raro.
\textbf{18.} Probabilidade a priori $\tfrac12$. Dois filhos saudáveis
têm probabilidade $\tfrac14$ se ela for portadora, 1 se não for:
probabilidade a posteriori $(\tfrac12\times
\tfrac14)/(\tfrac12\times\tfrac14 + \tfrac12) = \tfrac15$.
\textbf{19.} O único X de um filho vem da mãe, e o Y não leva nenhum
desses genes, de modo que cada filho exibe um gameta materno sem
mascaramento: o cruzamento é seu próprio cruzamento-teste.
\textbf{20.} Parentais $+++$ (853), $ywm$ (833); crossing-overs duplos
$+w+$ (2), $y+m$ (2); $w$ mudou: ordem $y$--$w$--$m$.
\textbf{21.} $r_{yw} = (24 + 26 + 2 + 2)/2000 = 0.027$; $r_{wm} = (128
+ 132 + 2 + 2)/2000 = 0.132$. Mapa: $y$ --- \qty{2.7}{cM} --- $w$ ---
\qty{13.2}{cM} --- $m$.
\textbf{22.} Esperados $0.027\times 0.132\times 2000 = 7.1$; observados
4: coincidência $0.56$, interferência $0.44$.
\textbf{23.} Recombinantes $y$--$m$: $y++$, $+wm$, $yw+$, $++m$: $(24 +
26 + 128 + 132)/2000 = 0.155$, contra $0.027 + 0.132 = 0.159$: os
quatro crossing-overs duplos restauram a combinação parental $y$--$m$
e passam despercebidos, $2\times 4/2000 = 0.004$.
\textbf{24.} As filhas recebem o X $y\,w\,m$ do pai e um X materno:
todas são heterozigotas e de \omterm{def:b2:meiosis-heredity:vocabulary}{fenótipo} selvagem, qualquer que seja o X
recombinante que carreguem; a recombinação é invisível.
\textbf{25.} $y$ --- \qty{2.7}{cM} --- $w$ --- \qty{13.2}{cM} --- $m$;
coeficiente de coincidência $0.56$, interferência $0.44$.
\end{solution}
